\section*{Bab \ref{ch:b2:structure-determination} --- Penentuan Struktur: NMR \texorpdfstring{$^{13}$C}{13C}, MS, dan IR Bersama-sama}

\begin{solution}{exo:b2:structure-determination:1}
1,2-Dimetilbenzena: 4 (\ce{CH3}; C1/C2; C3/C6; C4/C5). 1,3-: 5 (\ce{CH3}; C1/C3; C2; C4/C6; C5).
1,4-: 3 (\ce{CH3}; C1/C4; C2/C3/C5/C6).
\end{solution}

\begin{solution}{exo:b2:structure-determination:2}
\ce{CH3CH(OH)CH2CH3}. DEPT-135: C1 (\ce{CH3}) ke atas, C2 (CH) ke atas, C3 (\ce{CH2}) ke bawah, C4 (\ce{CH3}) ke atas.
DEPT-90: hanya C2.
\end{solution}

\begin{solution}{exo:b2:structure-determination:3}
209: karbonil keton; 170: karbonil ester (atau asam); 129: karbon \omterm{def:b2:huckel:aromatic}{aromatik};
64: karbon yang terikat pada oksigen; 14: \ce{CH3} alkil.
\end{solution}

\begin{solution}{exo:b2:structure-determination:4}
$(2 \times 8 + 2 - 8)/2 = 5$: sebuah cincin benzena (4) dan sebuah \ce{C=O}
(1). Misalnya metil benzoat,
\ce{C6H5COOCH3}, atau asam 4-metilbenzoat.
\end{solution}

\begin{solution}{exo:b2:structure-determination:5}
Satu ketidakjenuhan, sebuah keton (\qty{1715}{cm^{-1}}, garis pada 209,
kuaterner); satu \ce{CH2} dan dua
\ce{CH3}, empat karbon yang berbeda: butan-2-on, \ce{CH3COCH2CH3}.
\end{solution}

\begin{solution}{exo:b2:structure-determination:6}
Karbon: pentan-2-on memperlihatkan lima garis, pentan-3-on yang simetris
tiga. Proton: pentan-2-on mempunyai singlet tiga proton (\ce{CH3CO});
pentan-3-on hanya kuartet dan triplet. \omterm{def:b2:mass-spec-atomic:spectrum}{Spektrum massa}: pentan-2-on mempunyai
\omterm{def:b2:mass-spec-atomic:spectrum}{puncak dasar} pada 43 (\ce{CH3CO+}) dan ion McLafferty pada 58; pentan-3-on
mempunyai \omterm{def:b2:mass-spec-atomic:spectrum}{puncak dasar} pada 57
(\ce{C2H5CO+}).
\end{solution}

\begin{solution}{exo:b2:structure-determination:7}
\ce{C4H8O2}, 88: pasangan kuartet--triplet adalah \ce{OCH2CH3} (\ce{CH2} pada
4.12 terikat pada oksigen), singletnya \ce{CH3CO}. Etil etanoat,
\ce{CH3COOCH2CH3}.
\end{solution}

\begin{solution}{exo:b2:structure-determination:8}
1,2-Diklorobenzena: tiga garis (C1/C2, C3/C6, C4/C5). 1,4-Diklorobenzena: dua
(C1/C4, keempat CH). DEPT-90 memperlihatkan dua garis CH untuk yang pertama,
satu untuk yang kedua.
\end{solution}

\begin{solution}{exo:b2:structure-determination:9}
C2 adalah pusat stereogenik. Mengganti salah satu proton C3 dengan deuterium
menghasilkan dua diastereoisomer, bukan enansiomer: kedua proton itu
diastereotopik, berada dalam lingkungan yang berbeda. Keduanya dapat
mempunyai pergeseran yang berbeda dan saling terkopel sekaligus terkopel
dengan tetangga-tetangganya, sehingga
\ce{CH2} memberikan multiplet yang rumit alih-alih kuintet sederhana.
\end{solution}

\begin{solution}{exo:b2:structure-determination:10}
Lima ketidakjenuhan; cincin monosubstitusi (tiga garis CH dan satu garis C
pada 137.0, lima H \omterm{def:b2:huckel:aromatic}{aromatik}); sebuah keton (200.8, C); gugus etil (\ce{CH2}
31.8, kuartet pada 2.98 di sebelah karbonil, \ce{CH3}
8.3). 1-Fenilpropan-1-on, \ce{C6H5COCH2CH3}. Konjugasi dengan cincin
mendelokalisasikan elektron-elektron
$\pi$ \ce{C=O} dan melemahkan ikatannya: pitanya turun di bawah \qty{1715}{cm^{-1}} keton jenuh.
% ledger: ir:CO-ketone
\end{solution}

\begin{solution}{exo:b2:structure-determination:11}
Asam pentanoat: pita \ce{O-H} yang sangat lebar dari 2500 sampai
\qty{3300}{cm^{-1}} dan sebuah proton di dekat 11--12 ppm. Metil butanoat:
singlet tiga proton dalam rentang proton pada karbon yang terikat pada
oksigen (3.3--4.5). Etil propanoat: dua kuartet, satu pada oksigen dan satu
di sebelah karbonil. Propil etanoat: singlet tiga proton di sebelah karbonil
(2.0--2.4) dan triplet pada oksigen.
% ledger: ir:OH-acid, nmrr:acid, nmrr:alcohol, nmrr:ketone
\end{solution}

\begin{solution}{exo:b2:structure-determination:12}
Kedua garis lenyap dalam \omterm{def:b2:structure-determination:dept}{DEPT}, sehingga keduanya kuaterner, dan \omterm{def:b2:structure-determination:dept}{DEPT} saja
tidak dapat membedakannya. Pergeserannya dapat: karbonil ester terletak di
dekat 167--171 dan karbon \omterm{def:b2:huckel:aromatic}{aromatik} di antara 128 dan 137 dalam bagan bab
ini. Karbon cincin pada 166.5 dan karbonil pada 130.6 akan sama-sama berada
jauh di luar kelasnya; penetapan yang benar adalah kebalikannya.
\end{solution}

\begin{solution}{pb:b2:structure-determination:1}
\textbf{1.} $3.1/31.7 \approx 9.8~\%$; $9.8/1.11 \approx 9$ karbon.
\textbf{2.} Massa genap: tidak ada nitrogen (atau dua).
\textbf{3.} $150 - 108 = 42$: \ce{H10O2} cocok (10 + 32): \ce{C9H10O2}.
\textbf{4.} Sepuluh karbon akan memberikan $M+1 \approx 11~\%$ dari $M$,
bukan 9.8~\%; dan IR memperlihatkan karbonil, yang memerlukan oksigen kedua
selain pita \ce{C-O}.
\textbf{5.} $(2 \times 9 + 2 - 10)/2 = 5$.
\textbf{6.} $108 + 10 \times 1.007825 + 2 \times 15.994915 \approx 150.0681$.
\textbf{7.} Regangan \ce{C=O} pada posisi ester jenuh (di dekat \qty{1735}{cm^{-1}}).
\textbf{8.} Regangan \ce{C-O} ester.
\textbf{9.} Ikatan \ce{C-H} \omterm{def:b2:huckel:aromatic}{aromatik} (atau alkena).
\textbf{10.} Gugus \ce{O-H}: tidak ada alkohol, tidak ada asam karboksilat.
\textbf{11.} Tidak: karbonil yang terkonjugasi dengan cincin menyerap pada
bilangan gelombang yang lebih rendah; di sini karbonil itu berada di posisi
ester jenuh, sehingga ada gugus tak terkonjugasi yang memisahkannya dari
cincin.
\textbf{12.} 170.70: \ce{C=O} ester; 136.14: karbon cincin yang membawa
substituen; 128.56 dan 128.24: CH \omterm{def:b2:huckel:aromatic}{aromatik}; 66.24: \ce{CH2} yang terikat
pada oksigen; 20.82: \ce{CH3}.
\textbf{13.} 66.24, yaitu \ce{CH2}.
\textbf{14.} 170.70 dan 136.14, kedua karbon tanpa hidrogen.
\textbf{15.} 128.56 dan 128.24.
\textbf{16.} Empat: karbon yang membawa substituen, kedua CH orto, kedua CH
meta, dan CH para. Tiga jenis CH yang diharapkan; dua di antaranya
mempunyai pergeseran yang terlalu berdekatan untuk terpisah pada medan ini,
dan berbagi satu garis.
\textbf{17.} Belum tentu: tinggi bukanlah jumlah (\cref{prop:b2:structure-determination:intensities}),
walaupun di sini garis itu mungkin memang memuat dua jenis CH yang tumpang-tindih.
\textbf{18.} 7.33: kelima proton \omterm{def:b2:huckel:aromatic}{aromatik}; 5.09: \ce{OCH2}; 2.06: \ce{CH3CO}.
\textbf{19.} Atom-atom tetangganya (karbon cincin dan oksigen untuk \ce{CH2},
karbon karbonil untuk \ce{CH3}) tidak membawa hidrogen.
\textbf{20.} Gugus itu terikat pada oksigen ester yang elektronegatif dan
pada cincin: keduanya mengurangi perisainya.
\textbf{21.} \ce{C6H5CH2OCOCH3}, benzil etanoat.
\textbf{22.} Dalam isomer itu \ce{CH2} tidak terikat pada oksigen, sehingga
tidak mungkin muncul pada 5.09, dan metilnya, yang terikat pada oksigen,
akan muncul dalam rentang 3.3--4.5, bukan pada 2.06 di sebelah karbonil.
% ledger: nmrr:alcohol, nmrr:ketone
\textbf{23.} 108: kehilangan 42 (ketena, \ce{CH2=C=O}) melalui penataan
ulang, yaitu kation radikal
\ce{C7H8O}; 91: \ce{C7H7+}, kation benzil; 43: \ce{CH3CO+}.
\textbf{24.} Tujuh jenis karbon (\ce{C=O}, karbon cincin yang tersubstitusi,
CH orto, meta, dan para,
\ce{CH2}, \ce{CH3}): \textbf{tujuh garis, satu di antaranya (\ce{CH2}) mengarah ke bawah} dalam DEPT-135.
\end{solution}
